% !TEX program = xelatex
\documentclass[a4paper,10pt]{article}
\usepackage[UTF8,fontset=none]{ctex}
\usepackage{xeCJK}
\setCJKmainfont{Songti SC}
\setCJKsansfont{STHeiti}
\setCJKmonofont{Songti SC}
\usepackage[margin=2.2cm]{geometry}
\usepackage{tikz}
\usetikzlibrary{shapes.geometric, arrows.meta, positioning, calc, fit}
\usepackage{amsmath, amssymb}
\usepackage{longtable}
\usepackage{array}
\usepackage{booktabs}
\usepackage{multirow}
\usepackage{enumitem}
\usepackage{fancyhdr}
\usepackage{titlesec}
\usepackage{hyperref}
\usepackage{xcolor}
\usepackage{tabularx}

\hypersetup{
  colorlinks=true,
  linkcolor=blue!60!black,
  urlcolor=blue!60!black,
  bookmarksnumbered=true
}

\pagestyle{fancy}
\fancyhf{}
\fancyhead[L]{\small predictBicycleState 函数设计文档}
\fancyhead[R]{\small 内部公开}
\fancyfoot[C]{\thepage}
\renewcommand{\headrulewidth}{0.4pt}

\setcounter{secnumdepth}{3}
\setcounter{tocdepth}{3}

\newcolumntype{L}[1]{>{\raggedright\arraybackslash}p{#1}}
\newcolumntype{C}[1]{>{\centering\arraybackslash}p{#1}}

% ============================================================
\begin{document}

% ==================== 封面表格 ====================
\thispagestyle{empty}

\begin{center}
{\CJKfamily{zhsong}\bfseries\Large 函数设计文档}
\end{center}

\vspace{8pt}

\begin{center}
\renewcommand{\arraystretch}{1.5}
\begin{tabular}{|L{3.2cm}|L{4.5cm}|L{3.2cm}|L{4.5cm}|}
\hline
\textbf{函数英文名} & predictBicycleState &
\textbf{函数中文名} & 自行车模型状态预测 \\
\hline
\textbf{函数类型} & 元件 (Element) &
\textbf{版本} & 2.0 \\
\hline
\textbf{作者} & 许庆 &
\textbf{日期} & 2026-06-26 \\
\hline
\textbf{联系方式} & \multicolumn{3}{l|}{18612426814~/~qingxu@tsinghua.edu.cn} \\
\hline
\textbf{密级} & 内部公开 &
\textbf{AI辅助} & Cursor Agent (Claude Opus 4.6) \\
\hline
\end{tabular}
\end{center}

\vspace{12pt}

\begin{center}
\renewcommand{\arraystretch}{1.4}
\begin{tabular}{|C{1cm}|C{1.8cm}|L{3.8cm}|C{1.5cm}|C{1.5cm}|C{3cm}|}
\hline
\multicolumn{6}{|c|}{\textbf{更改记录}} \\
\hline
\textbf{版本} & \textbf{日期} & \textbf{更改说明} & \textbf{作者} & \textbf{审核} & \textbf{辅助AI} \\
\hline
1.0 & 2026-06-20 & 初始版本 & 许庆 & --- & --- \\
\hline
2.0 & 2026-06-26 & 统一接口为结构体传参 & 许庆 & --- & Claude Opus 4.6 \\
\hline
\end{tabular}
\end{center}

\newpage

% ==================== 目录 ====================
\tableofcontents
\newpage

% ============================================================
%  1. 函数需求
% ============================================================
\section{函数需求}

\subsection{功能概述}

\texttt{predictBicycleState} 基于运动学自行车模型，使用欧拉前向离散化方法
预测下一时刻的车辆状态。给定当前车辆状态（位置、航向、速度）和控制输入
（前轮转向角、加速度），以及时间步长和轴距，输出预测的下一时刻车辆状态。

\subsection{输入/输出/返回值}

\renewcommand{\arraystretch}{1.3}
\begin{longtable}{|L{2.2cm}|L{3cm}|C{2.5cm}|L{7cm}|}
\hline
\textbf{方向} & \textbf{名称} & \textbf{类型} & \textbf{说明} \\
\hline
\endfirsthead
\hline
\textbf{方向} & \textbf{名称} & \textbf{类型} & \textbf{说明} \\
\hline
\endhead
输入 & \texttt{state} & \texttt{VehicleState\&} & 当前车辆状态 $(x, y, \psi, v)$ \\
\hline
输入 & \texttt{control} & \texttt{ControlInput\&} & 控制输入 $(\delta, a)$ \\
\hline
输入 & \texttt{dt} & \texttt{double} & 时间步长 (s) \\
\hline
输入 & \texttt{wheelbase} & \texttt{double} & 前后轴距 $L$ (m) \\
\hline
返回值 & --- & \texttt{VehicleState} & 预测的下一时刻车辆状态 \\
\hline
\end{longtable}

\subsection{功能需求}

\begin{enumerate}[label=FR-\arabic*]
  \item 根据运动学自行车模型的离散方程计算下一时刻的位置 $(x_{k+1}, y_{k+1})$。
  \item 计算下一时刻的航向角 $\psi_{k+1}$。
  \item 计算下一时刻的纵向速度 $v_{k+1}$。
  \item 所有中间量（$\cos\psi$, $\sin\psi$, $\tan\delta$）仅计算一次。
\end{enumerate}

\subsection{非功能需求}

\begin{enumerate}[label=NFR-\arabic*]
  \item 不使用动态内存分配。
  \item 仅依赖 \texttt{<cmath>} 标准库。
  \item 纯函数：无内部状态，相同输入保证相同输出。
\end{enumerate}

% ============================================================
%  2. 算法设计
% ============================================================
\section{算法设计}

\subsection{运动学自行车模型}

采用前轮转向运动学自行车模型（Kinematic Bicycle Model），
将车辆简化为两轮模型，后轴中心为参考点。

状态向量 $\mathbf{s} = [x, y, \psi, v]^{\!\top}$，控制向量 $\mathbf{u} = [\delta, a]^{\!\top}$。
连续时间运动学方程为：
\begin{equation}
  \begin{cases}
    \dot{x} = v \cos\psi \\[4pt]
    \dot{y} = v \sin\psi \\[4pt]
    \dot{\psi} = \dfrac{v \tan\delta}{L} \\[6pt]
    \dot{v} = a
  \end{cases}
\end{equation}
其中：
\begin{itemize}
  \item $x, y$：后轴中心在全局坐标系中的位置 (m)
  \item $\psi$：航向角 (rad)，逆时针为正
  \item $v$：纵向速度 (m/s)
  \item $\delta$：前轮转向角 (rad)，左转为正
  \item $a$：纵向加速度 (m/s\textsuperscript{2})
  \item $L$：前后轴距 (m)
\end{itemize}

\subsection{欧拉前向离散化}

采用欧拉前向差分（显式欧拉法），以时间步长 $\Delta t$ 进行离散化：
\begin{equation}\label{eq:discrete}
  \boxed{
  \begin{aligned}
    x_{k+1}   &= x_k + v_k \cos\psi_k \cdot \Delta t \\
    y_{k+1}   &= y_k + v_k \sin\psi_k \cdot \Delta t \\
    \psi_{k+1} &= \psi_k + \frac{v_k \tan\delta_k}{L} \cdot \Delta t \\
    v_{k+1}   &= v_k + a_k \cdot \Delta t
  \end{aligned}
  }
\end{equation}

航向角变化率（横摆角速度）为：
\begin{equation}
  \dot{\psi} = \frac{v \tan\delta}{L}
\end{equation}

当 $\delta \to \pm\frac{\pi}{2}$ 时，$\tan\delta \to \pm\infty$，
导致横摆角速度趋于无穷大。实际使用中上层控制器负责将转向角限制在合理范围内。

% ============================================================
%  3. 程序流程图
% ============================================================
\section{程序流程图}

\begin{center}
\begin{tikzpicture}[
  node distance=0.9cm,
  startstop/.style={rectangle, rounded corners, minimum width=3.5cm, minimum height=0.7cm,
    text centered, draw=black, fill=blue!8, font=\small},
  process/.style={rectangle, minimum width=5.5cm, minimum height=0.7cm,
    text centered, draw=black, fill=orange!8, font=\small},
  io/.style={trapezium, trapezium left angle=70, trapezium right angle=110,
    minimum width=3.5cm, minimum height=0.7cm,
    text centered, draw=black, fill=yellow!8, font=\small},
  arrow/.style={-{Stealth[length=2.5mm]}, thick}
]

\node (start) [startstop] {开始};
\node (input) [io, below=of start] {读入 state, control, dt, wheelbase};
\node (trig) [process, below=of input] {计算 $\cos\psi$, $\sin\psi$, $\tan\delta$};
\node (yawrate) [process, below=of trig] {计算 yawRate $= v \cdot \tan\delta / L$};
\node (xnext) [process, below=of yawrate] {$x_{k+1} = x + v \cdot \cos\psi \cdot \Delta t$};
\node (ynext) [process, below=of xnext] {$y_{k+1} = y + v \cdot \sin\psi \cdot \Delta t$};
\node (yawnext) [process, below=of ynext] {$\psi_{k+1} = \psi + \text{yawRate} \cdot \Delta t$};
\node (vnext) [process, below=of yawnext] {$v_{k+1} = v + a \cdot \Delta t$};
\node (output) [io, below=of vnext] {输出 nextState$(x_{k+1}, y_{k+1}, \psi_{k+1}, v_{k+1})$};
\node (stop) [startstop, below=of output] {结束};

\draw [arrow] (start) -- (input);
\draw [arrow] (input) -- (trig);
\draw [arrow] (trig) -- (yawrate);
\draw [arrow] (yawrate) -- (xnext);
\draw [arrow] (xnext) -- (ynext);
\draw [arrow] (ynext) -- (yawnext);
\draw [arrow] (yawnext) -- (vnext);
\draw [arrow] (vnext) -- (output);
\draw [arrow] (output) -- (stop);

\end{tikzpicture}
\end{center}

% ============================================================
%  4. 函数接口与输入输出模块图
% ============================================================
\section{函数接口与输入输出模块图}

\subsection{函数接口}

\begin{verbatim}
VehicleState predictBicycleState(const VehicleState& state,
                                  const ControlInput& control,
                                  double dt,
                                  double wheelbase);
\end{verbatim}

\subsection{输入输出模块图}

\begin{center}
\begin{tikzpicture}[
  core/.style={rectangle, draw=black, very thick, minimum width=5cm, minimum height=2cm,
    text centered, fill=orange!10, font=\normalsize\bfseries, rounded corners=3pt},
  iobox/.style={rectangle, draw=black, thick, minimum width=3.2cm, minimum height=0.6cm,
    text centered, fill=yellow!10, font=\small},
  arrow/.style={-{Stealth[length=2.5mm]}, thick}
]

\node (core) [core] {predictBicycleState};

\node (in1) [iobox, left=3cm of core, yshift=1.2cm] {state (VehicleState)};
\node (in2) [iobox, left=3cm of core, yshift=0.4cm] {control (ControlInput)};
\node (in3) [iobox, left=3cm of core, yshift=-0.4cm] {dt (double)};
\node (in4) [iobox, left=3cm of core, yshift=-1.2cm] {wheelbase (double)};

\node (out1) [iobox, right=3cm of core] {VehicleState (nextState)};

\draw [arrow] (in1) -- (core);
\draw [arrow] (in2) -- (core);
\draw [arrow] (in3) -- (core);
\draw [arrow] (in4) -- (core);
\draw [arrow] (core) -- (out1);

\node [above=0.1cm of in1, font=\small\bfseries] {输入 (Input)};
\node [above=0.1cm of out1, font=\small\bfseries] {输出 (Output)};

\end{tikzpicture}
\end{center}

% ============================================================
%  5. 函数诸元表
% ============================================================
\section{函数诸元表}

\renewcommand{\arraystretch}{1.3}
\begin{longtable}{|L{3.8cm}|L{11.5cm}|}
\hline
\textbf{属性} & \textbf{说明} \\
\hline
\endfirsthead
\hline
\textbf{属性} & \textbf{说明} \\
\hline
\endhead
函数英文名 & \texttt{predictBicycleState} \\
\hline
函数中文名 & 自行车模型状态预测 \\
\hline
函数类型 & 元件 (Element) \\
\hline
版本 & 2.0 \\
\hline
源文件 & \texttt{src/cpp/predictBicycleState/predictBicycleState.cpp} \\
\hline
头文件 & \texttt{include/cpp/predictBicycleState/predictBicycleState.hpp} \\
\hline
命名空间 & \texttt{waypoint\_follow} \\
\hline
输入数量 & 4（state, control, dt, wheelbase） \\
\hline
输出数量 & 1（VehicleState 返回值） \\
\hline
参数数量 & 0 \\
\hline
状态数量 & 0（纯函数） \\
\hline
下级函数数量 & 0 \\
\hline
动态内存分配 & 无 \\
\hline
外部依赖 & \texttt{<cmath>}（\texttt{std::cos}, \texttt{std::sin}, \texttt{std::tan}） \\
\hline
算法类别 & 运动学模型 / 欧拉离散化 \\
\hline
时间复杂度 & $O(1)$ \\
\hline
\end{longtable}

% ============================================================
%  6. 函数架构图
% ============================================================
\section{函数架构图}

\begin{center}
\begin{tikzpicture}[
  every node/.style={rectangle, draw, rounded corners=2pt, thick,
    minimum height=0.65cm, text centered, font=\footnotesize,
    fill=blue!5}
]

\node (main) {predictBicycleState};
\node [below=0.6cm of main, draw=none, fill=none, font=\small\itshape] {无下级函数调用};

\end{tikzpicture}
\end{center}

% ============================================================
%  7. 下级函数需求
% ============================================================
\section{下级函数需求}

本函数为叶子元件，无下级函数调用。

% ============================================================
%  8. 数据结构定义
% ============================================================
\section{数据结构定义}

以下数据结构定义于 \texttt{waypointFollowTypes.hpp}，命名空间 \texttt{waypoint\_follow}。

\subsection{VehicleState~---~车辆运动状态}

\begin{longtable}{|L{3.5cm}|C{2cm}|L{9.5cm}|}
\hline
\textbf{成员} & \textbf{类型} & \textbf{说明} \\
\hline
\texttt{x} & double & 纵向位置 (m) \\
\hline
\texttt{y} & double & 横向位置 (m) \\
\hline
\texttt{yaw} & double & 航向角 (rad)，逆时针为正 \\
\hline
\texttt{v} & double & 纵向速度 (m/s) \\
\hline
\end{longtable}

\subsection{ControlInput~---~车辆控制输入量}

\begin{longtable}{|L{3.5cm}|C{2cm}|L{9.5cm}|}
\hline
\textbf{成员} & \textbf{类型} & \textbf{说明} \\
\hline
\texttt{steeringAngle} & double & 前轮转向角 (rad)，左转为正 \\
\hline
\texttt{acceleration} & double & 纵向加速度 (m/s\textsuperscript{2})，加速为正 \\
\hline
\end{longtable}

% ============================================================
%  9. 测试用例
% ============================================================
\section{测试用例}

\subsection{正常测试用例}

{\small
\begin{longtable}{|C{0.6cm}|L{1.8cm}|L{5.0cm}|L{3.8cm}|L{3.5cm}|}
\hline
\textbf{编号} & \textbf{场景} & \textbf{输入} & \textbf{预期输出} & \textbf{验证准则} \\
\hline
\endfirsthead
\hline
\textbf{编号} & \textbf{场景} & \textbf{输入} & \textbf{预期输出} & \textbf{验证准则} \\
\hline
\endhead
N-1 &
直线行驶 &
$\mathbf{s}{=}(0,0,0,10)$, $\delta{=}0$, $a{=}0$, $\Delta t{=}0.1$, $L{=}2.85$ &
$(1.0, 0, 0, 10)$ &
位置沿 $x$ 轴前进，$\psi$, $v$ 不变 \\
\hline
N-2 &
恒定左转弯 &
$\mathbf{s}{=}(0,0,0,10)$, $\delta{=}0.1$, $a{=}0$, $\Delta t{=}0.1$, $L{=}2.85$ &
$(1.0, 0, 0.035, 10)$ &
$\psi_{k+1} > 0$（左转），$v$ 不变 \\
\hline
N-3 &
静止加速 &
$\mathbf{s}{=}(0,0,0,0)$, $\delta{=}0$, $a{=}2$, $\Delta t{=}0.1$, $L{=}2.85$ &
$(0, 0, 0, 0.2)$ &
位置不变，$v_{k+1}{=}0.2$ \\
\hline
N-4 &
制动减速 &
$\mathbf{s}{=}(0,0,0,20)$, $\delta{=}0$, $a{=}{-}3$, $\Delta t{=}0.1$, $L{=}2.85$ &
$(2.0, 0, 0, 19.7)$ &
$v_{k+1} < v_k$, $x_{k+1} > 0$ \\
\hline
N-5 &
斜向行驶 &
$\mathbf{s}{=}(0,0,\pi/4,10)$, $\delta{=}0.05$, $a{=}0$, $\Delta t{=}0.1$, $L{=}2.85$ &
$x_{k+1} \approx 0.707$, $y_{k+1} \approx 0.707$ &
沿 $45^{\circ}$ 方向前进 \\
\hline
\end{longtable}
}

\vspace{6pt}
{\small
\begin{longtable}{|L{2.8cm}|L{12.5cm}|}
\hline
\multicolumn{2}{|c|}{\textbf{正常测试用例符号说明}} \\
\hline
\textbf{符号} & \textbf{含义} \\
\hline
$\mathbf{s}{=}(x,y,\psi,v)$ & 车辆状态 VehicleState \\
\hline
$\delta$ & 前轮转向角 steeringAngle (rad) \\
\hline
$a$ & 纵向加速度 acceleration (m/s\textsuperscript{2}) \\
\hline
$\Delta t$ & 时间步长 dt (s) \\
\hline
$L$ & 前后轴距 wheelbase (m) \\
\hline
$\psi$ & 航向角 yaw (rad) \\
\hline
$v$ & 纵向速度 (m/s) \\
\hline
\end{longtable}
}

\subsection{异常测试用例}

{\small
\begin{longtable}{|C{0.6cm}|L{1.8cm}|L{4.8cm}|L{3.8cm}|L{3.5cm}|}
\hline
\textbf{编号} & \textbf{场景} & \textbf{输入} & \textbf{预期输出/行为} & \textbf{验证准则} \\
\hline
\endfirsthead
\hline
\textbf{编号} & \textbf{场景} & \textbf{输入} & \textbf{预期输出/行为} & \textbf{验证准则} \\
\hline
\endhead
A-1 &
零速+转向 &
$v{=}0$, $\delta{=}0.3$, $\Delta t{=}0.1$ &
位置不变, $\psi$ 不变, $v{=}0$ &
yawRate $= 0 \cdot \tan\delta / L = 0$ \\
\hline
A-2 &
速度为 NaN &
$v{=}\text{NaN}$, 其余正常 &
不崩溃; 输出各分量可能为 NaN &
函数正常返回 \\
\hline
A-3 &
加速度 $+\infty$ &
$a{=}+\infty$, $v{=}10$, $\Delta t{=}0.1$ &
不崩溃; $v_{k+1}{=}+\infty$ &
无浮点异常 \\
\hline
A-4 &
$\delta{=}\pi/2$ &
$\delta{=}\pi/2$, $v{=}10$, $L{=}2.85$ &
$\tan(\pi/2) \to \infty$; $\dot\psi \to \infty$ &
yawRate 为极大值或 $\pm\infty$ \\
\hline
A-5 &
轴距为零 &
$L{=}0$, $v{=}10$, $\delta{=}0.1$ &
除零: $v\tan\delta/0 \to \pm\infty$ &
$\psi_{k+1}$ 为 $\pm\infty$ 或 NaN \\
\hline
A-6 &
负轴距 &
$L{=}{-}2.85$, $v{=}10$, $\delta{=}0.1$ &
航向变化符号反转 &
$\dot\psi$ 符号与正轴距相反 \\
\hline
A-7 &
航向角 NaN &
$\psi{=}\text{NaN}$, 其余正常 &
不崩溃; $\cos/\sin$ 输出 NaN &
位置增量为 NaN \\
\hline
A-8 &
极大速度 &
$v{=}10^{15}$, $\Delta t{=}0.1$ &
位移 $= 10^{14}$ (极大) &
不溢出为 $\pm\infty$ \\
\hline
A-9 &
负时间步长 &
$\Delta t{=}{-}0.1$, $v{=}10$, $\delta{=}0$ &
$x_{k+1} < x_k$（时间反向） &
状态向反方向演化 \\
\hline
A-10 &
零时间步长 &
$\Delta t{=}0$, 其余正常 &
状态完全不变 &
$\mathbf{s}_{k+1} = \mathbf{s}_k$ \\
\hline
\end{longtable}
}

\vspace{6pt}
{\small
\begin{longtable}{|L{2.8cm}|L{12.5cm}|}
\hline
\multicolumn{2}{|c|}{\textbf{异常测试用例符号说明}} \\
\hline
\textbf{符号} & \textbf{含义} \\
\hline
$v$ & 纵向速度 (m/s) \\
\hline
$\delta$ & 前轮转向角 steeringAngle (rad) \\
\hline
$a$ & 纵向加速度 acceleration (m/s\textsuperscript{2}) \\
\hline
$\psi$ & 航向角 yaw (rad) \\
\hline
$L$ & 前后轴距 wheelbase (m) \\
\hline
$\Delta t$ & 时间步长 dt (s) \\
\hline
$\dot\psi$ & 横摆角速度 yawRate (rad/s) \\
\hline
NaN & IEEE 754 非数值 (Not a Number) \\
\hline
$+\infty$ & IEEE 754 正无穷大 \\
\hline
\end{longtable}
}

\end{document}
